% Folder 93, batch 07 — archivists' pages 121–140.
% Pass: Opus 5.5 (claude-opus-5-5), 2026-10-10 — first pass, unchecked against the pages by a human.
%
% Pages 124 and 140 are blank separator sheets and are not transcribed.
\documentclass[11pt,a4paper]{article}
\input{../preamble/grothendieck}

\folder{93}
\batch{7}
\pages{121}{140}
\foldertitle{Connexions de Gauss-Manin et équations de Picard-Fuchs. Opération de Cartier : copies de tapuscrit annoté (s.d.), tiré à part (1981), notes manuscrites (s.d.).}
\dating{[1972]-1981}
\watermark{Édition de démonstration}

\begin{document}

\page{121}
\note{la page commence au milieu d'une phrase : l'argument vient de la page qui précède le lot}

des \add{$\mathrm{Conn}_{X/S}$ et des} $H^{1}\mathcal{P}_{X/S} \longrightarrow \mathrm{Conn}_{X/S}$ s'étendant à $G_0(\ill{})$ \uncertain{et de}
\[
H^{1}\mathcal{P}_{X/S} \longleftrightarrow \mathrm{Conn}^{0}_{X/S}
\]
les premiers se déduisant du second de façon évidente, compte tenu de l'isom.\ can.
\[
\mathrm{Conn}^{0}_{X/S} \longrightarrow \mathrm{Conn}_{X/S}
\]
compatible avec l'\uncertain{iso}
\[
\underline{\omega}_{X/S}^{\otimes 2} \xrightarrow{\ 12\,\mathrm{id}\ } \underline{\omega}_{X/S}^{\otimes 2}
\]
sur les \uncertain{champs} d'opérateurs. (On \add{écrira} \struck{\ill{}} \uncertain{parfois},
\[
c_{0}^{X}(\ill{}) = \tfrac{1}{12}\, c^{X}(\ill{}) \quad \ldots\,)
\]

On peut espérer que si $X$ est de car.\ $p > 0$, et si $\mathcal{O}'_X$ est le sous-faisceau \add{d'anneaux \uncertain{des fonctions}} de $\mathcal{O}_X$ \struck{défini par} des $f$ tels que $D^{(i)}f = 0$ ($i \leqslant 3$), \struck{\ill{}} \uncertain{les} conditions
\[
H_{\mathcal{X}}(\mathcal{P}/S) \backslash \mathcal{P}_{X/S} \longleftrightarrow \mathrm{Conn}^{0}_{X/S}
\]
\note{l'indice de $H$ surcharge un premier symbole ; lecture douteuse}

sont \uncertain{injectifs}, i.e.\ \uncertain{que} \ill{}

\marginal{$\mathcal{O}'_{X/S}$ est $\mathcal{O}^{(p)}_{X/S}$ si $p \neq 2, 3$, \uncertain{sinon} $\mathcal{O}^{(p^{2})}_{X/S}$}
\marginal{OK}

\emph{Proposition} \struck{\ill{}} $A$ anneau (local, plus généralement avec $\operatorname{Pic}(A) = 0$) de car.\ $p > 0$, $\alpha \in A$, $B = A[T]/(T^{p} - \alpha)$, $t \in B$ la classe \struck{de} \add{de} $T$ dans $B$ (donc $t^{p} = \alpha$), \struck{\ill{}} $s \in B$, donc
\[
s = a_0 + a_1 t + \cdots + a_{p-1} t^{p-1}, \qquad a_i \in A \ (0 \leqslant i \leqslant p-1),
\]
\[
D_t s = (ds)/(dt), \qquad D^{(i)}_t s = \tfrac{1}{i!}(D_t)^{i}s \quad\text{si } 1 \leqslant i \leqslant 3,
\]
\struck{$D^{(j)}_t s = \sum_{j \leqslant i \leqslant p-1} a_i \binom{i}{j} t^{i-j}$, et}

\marginal{cas $p \neq 2, 3$}
\marginal{\uncertain{\ill{} \ill{}} (\ill{})}

Posons \struck{Alors}
\[
c_0(s) = -\frac{D^{(3)}_t s}{D^{(1)}_t s} + \Bigl(\frac{D^{(2)}_t s}{D^{(1)}_t s}\Bigr)^{2} \quad \text{si } D^{(1)}_t s \in B^{*}.
\]
Supposons $c_0(s) = 0$ i.e.
\[
(D^{(2)}s)^{2} = (D^{(1)}s)(D^{(3)}s).
\]
Alors $\exists\, a, b, c, d \in A$ (\uncertain{uniques} \ill{} \ill{} \ill{} \add{\uncertain{scalaire}}) \add{avec $ad - bc = 1$}, tels que
\[
s = \frac{at + b}{ct + d} \quad\text{i.e.}\quad s(ct + d) = at + b
\]
\note{les exposants (2) et (3) se distinguent mal dans ces deux formules ; on les lit d'après les calculs de la page 122, où $D^{(3)}s\,D^{(1)}s$ et $(D^{(2)}s)^{2}$ sont comparés}

\page{122}
OPS \struck{\ill{}} pour les éléments \struck{\ill{}} $\in A^{p}$, puis $\alpha \neq 0$ \struck{\ill{}}

\struck{donc} $t \in B$, donc $D^{(1)}s = a_1 + a_2 t + \cdots \equiv a_1 \ (tB)$\note{lecture de « $(tB)$ » douteuse}

donc $D^{(1)}s \in B^{*} \Longleftrightarrow a_1 \in A^{*}$. OPS, quitte à remplacer $s$ par $a_1^{-1}(s - a_0)$, que $s = t + a_2 t^{2} + \cdots + a_{p-1}t^{p-1}$, puis, quitte à faire encore un chgt de variable
\[
s' = \frac{s}{1 + a_2 s} = \sum_{n \geqslant 0} (-1)^{n} a_2^{n} s^{n+1}
= s - a_2 s^{2} + \cdots \equiv t \bmod t^{3}
\]
\note{entre la fraction et la somme, une seconde fraction $\frac{1}{1+a_2 s}$ est biffée ; devant le $\sum$ un $s$ semble biffé}

donc $s$ de la forme
\[
s = t + \text{termes d'ordre} \geqslant 3 .
\]

\struck{\ill{}} À prouver que, \uncertain{sous une condition}, la condition
\[
c_0(s) = 0 \Longrightarrow s = t .
\]
\struck{Si \ill{}} \ill{} par si :
\[
s = t + a t^{d} + \cdots \quad (\text{termes d'ordre} \geqslant d+1),
\]
$3 \leqslant d \leqslant p-1$ \struck{$d \geqslant 3$}, alors $a = 0$.
\begin{align*}
D^{(3)}s &\equiv a\binom{d}{3} t^{d-3} \bmod (t^{d-2})\\
D^{(2)}s &\equiv a\binom{d}{2} t^{d-2} \bmod (t^{d-1}) \text{ et a fortiori mod } t^{d-2}\\
D^{(1)}s &\equiv 1 \bmod (t^{d-1}) \text{ et a fortiori mod } t^{d-2}\\
D^{(3)}s\, D^{(1)}s &\equiv a\binom{d}{3} t^{d-3} \bmod t^{d-2}\\
(D^{(2)}s)^{2} &\equiv 0 \bmod t^{d-2}
\end{align*}
\note{à la ligne de $D^{(1)}s$, un terme $+\,a d\, t^{d-1}$ est biffé ; à celle de $D^{(3)}s\,D^{(1)}s$, un terme $+\,a^{2} d \binom{d}{3} t^{2d-4}$ est biffé ; à celle de $(D^{(2)}s)^{2}$, le $0$ surcharge un $a^{2}$}

donc
\[
a\binom{d}{3} t^{d-3} \equiv 0 \bmod t^{d-2}
\]
donc $a\binom{d}{3} = 0$ dans $A$, i.e.
\[
\frac{d(d-1)(d-2)}{6}\, a = 0 \text{ dans } A .
\]
Si $6$ \uncertain{inv.} dans $A$, et aussi $d(d-1)(d-2)$, \uncertain{car} $3 \leqslant d \leqslant p-1$ et $A$ de car. $p > 3$, donc $a = 0$ OK.

\page{123}
\emph{Cas car. 2 ou 3}. Ici il faut supposer $B = A[T]/(T^{p^{2}} - \alpha)$, dans l'\uncertain{unicité} (\uncertain{modulo vérification} de la \struck{\ill{}} \add{\uncertain{définition}} des $D^{(i)}_{t}$ …) du résultat sur \uncertain{formes} d'\ill{} \uncertain{connues}, et la réduction au cas $\alpha = 0$, $s \equiv t \ (t^{3}B)$, est comme dessus. Il faut prouver que $s = t$, donc si $s = t + a t^{d}$ ($3 \leqslant d \leqslant$ \struck{\ill{}} $p^{2}-1$) alors $a = 0$. On trouve comme \struck{\ill{}}
\[
a\binom{d}{3} = 0 \quad\text{i.e.}\quad a\, c_{3,d-3} = 0,
\]
\note{sous $c_{3,d-3}$, une flèche renvoie à « coeff. binômial »}

d'autre part $c_{3,d-3} \in \mathbb{F}_p^{*}$ ssi. \struck{$d \geqslant d-3$}
\[
\text{pour}\quad
\begin{aligned}
d-3 &= d_0 + d_1 p\\
3 &= d'_0 + d'_1 p
\end{aligned}
\qquad \Bigl|\ 0 \leqslant d_0, d_1, d'_0, d'_1 \leqslant p-1
\]
on a $d_0 + d'_0 \leqslant p-1$, $d_1 + d'_1 \leqslant p-1$.

1) Cas car. $p = 2$, \struck{donc $d - 3 = \ill{}\,(d_0 - 1) + (d_1 - 1)p$ \ill{}} \struck{On $d_0, d_1 \geqslant 1$ \ill{} \ill{}} Cas : signifie \struck{\ill{}} pour $d = 3$ \struck{\ill{}}
\[
c_{3,3} = 1 \Longrightarrow a = 0 \qquad \text{OK}
\]
\struck{\ill{}}

2) \emph{Cas $p = 3$.} \struck{\ill{}} Si $d_1 \neq 0$ i.e. $d_1 \geqslant 1$,
\[
d - 3 = d - p = d_0 + (d_1 - 1)p, \quad\text{et}
\]
\[
c_{d,d-3} \equiv c_{d_0,d_0}\, c_{d_1,d_1-1} \quad \text{dans } \mathbb{F}_p,
\]
\[
\in \mathbb{F}_p^{*} \text{ ssi }
\begin{aligned}
2d_0 &\leqslant p-1 = 2\\
2d_1 - 1 &\leqslant p-1 = 2
\end{aligned}
\]
i.e. $d_1 = 1$, $d_0 = 0$ ou $1$, dans \uncertain{ce qui ici dépend} …
(a) $d_1 = 0$ ; \struck{(a)} $d_1 = 2$. b) $d_1 = 1$, $d_0 = 2$. \struck{\ill{}} Le cas $d_1 = 0$ \ill{} on fait \uncertain{facilement}, \ill{} alors $3 \leqslant d \leqslant p-1$ dans $c_{d,d-3}$
\note{tout ce cas 2) est barré de trois longues diagonales ; « (a) $d_1 = 0$ » est cerclé}

\marginal{$d - 3 = d_0 + d_1 p$, $0 \leqslant d_0, d_1 \leqslant 2$ ; $0 \leqslant d_1 \leqslant 1$ ; donc $d_1 \leqslant 1$ (\uncertain{de même} $0 \leqslant d-3 \leqslant 5$ car $3 \leqslant d \leqslant p^{2}-1 = 8$). $3 = 0 + 1\cdot p$. La condition $c_{3,d-3} \in \mathbb{F}_p^{*}$ signifie ici $d_1 + 1 \leqslant p-1$ i.e. $d_1 + 1 \leqslant 2$ i.e. $d_1 \leqslant 1$, qui \uncertain{est} vérifié. On gagne encore !}
\note{cette colonne de droite, écrite après coup en haut à droite de la feuille, porte sur le cas $p = 3$ ; elle est suivie d'un grand paraphe}

\page{125}
Soit $P$ une droite projective sur $\mathbb{C}$. Un ouvert non vide $D \subset P^{\mathrm{an}}$ est dit « \emph{homogène} » si le \add{sous-}groupe $G$ de $\operatorname{Aut}(P) \xrightarrow{\sim} \operatorname{Aut}(P^{\mathrm{an}})$ ($\simeq GP(1)_{\mathbb{C}}$) agit homogènement sur $D$. Si $D$ est homogène et $g \in \operatorname{Aut}(P)$, $g(D)$ est homogène et a comme groupe de stabilité $\operatorname{int}(g).G$.

Le « théorème d'uniformisation » nous donne

\emph{Th.\ (classiques) 1}

a) Il y a exactement trois classes (à conj.\ près \uncertain{par} $\operatorname{Aut}(P)$) de domaines homogènes dans $P$, savoir
\begin{itemize}
\item[1°)] $D = P$
\item[2°)] $D = P - \{x\}$, $x \in P$
\item[3°)] $D$ un disque ouvert [pour une métrique (\add{riemannienne}) invariante sur $X$ — \uncertain{les} métriques \add{invariantes} sont toutes multiples l'une d'elles, et les disques ouverts \add{dans} pour une \add{telle} métrique sont indépendants du choix de celle-ci]
\end{itemize}

b) Pour $D$ domaine homogène dans $X$, $G \subset \operatorname{Aut}(P)$ son stabilisateur, l'application $G \to \operatorname{Aut}(D)$ est un iso. De plus, $D$ est connexe et \add{simplement} connexe.

c) Le foncteur $(P, D) \longmapsto D$, de la catégorie \add{(pour iso)} des couples $(P, D)$, $P$ droite projective et $D$ domaine homogène de $P$, dans la catégorie \add{(pour iso)} des \struck{\ill{}} espaces $\mathbb{C}$-analytiques réguliers de dim~1 connexes et simplement connexes, est une équivalence de catégories.

\emph{Cor.\ \uncertain{Naïve}} Les groupes $G$ de stabilité d'un domaine homogène $D$ de $P$ sont respectivement
\begin{itemize}
\item[1)] $\operatorname{Aut}(P)$ ($\simeq GP(1)_{\mathbb{C}}$)
\item[2)] $\operatorname{Aut}(\mathbb{E}^{1})$ ($\simeq \operatorname{Aff}(1, \mathbb{C})$)
\item[3)] $\simeq GP(1, \mathbb{R}) \simeq Gl(2, \mathbb{R})/\mathbb{R}^{*} \simeq Sl(2, \mathbb{R})/\mu_2(\mathbb{R})$
\end{itemize}
(On notera les affirmations 1°) et 2°) d'une \ill{} \ill{}, où on trouve (\struck{\ill{}} groupes analytiques \uncertain{complexes}) \ill{} une \uncertain{agissant} sur $\mathbb{C}$, et le cas 3°)

\page{126}
où on trouve \struck{un} \add{le groupe des pts réels $GP(1)(\mathbb{R})$} \struck{groupe algébrique réel $GP(1)_{\mathbb{R}}$} — qui en fait est une forme réelle de $GP(1)_{\mathbb{C}}$, \struck{\ill{}} \add{dans 3°)} \struck{\ill{}} \add{on trouve} que la \struck{\ill{}} \add{correspondance} \struck{forme} biunivoque entre
\begin{itemize}
\item[$\alpha$)] choix d'une \struck{forme réelle} \add{descente du corps de base de $\mathbb{C}$ à $\mathbb{R}$} dans le groupe alg.\ $\operatorname{Aut}_{\mathbb{C}} P$
\item[$\beta$)] choix d'une descente du corps de base de $\mathbb{C}$ à $\mathbb{R}$ dans l'algèbre de Lie complexe
\[
\operatorname{Lie}\operatorname{Aut}_{\mathbb{C}} P = H^{0}(P, \mathcal{V}_{P/\mathbb{C}}) \simeq H^{0}(P, \Omega^{-1}_{P/\mathbb{C}})
\]
\item[$\gamma$)] choix d'un domaine homogène du type 3°) (i.e.\ non algébrique) dans $P$, à polarité près (la polaire de $D$ est $P - \overline{D}$)
\end{itemize}
\marginal{$\delta$) choix d'une \ill{} \uncertain{géodésique} \ill{} ($= \ill{}$) dans $P$}

La relation entre $\alpha$) et $\beta$) claire, celle entre $\alpha$) et $\gamma$) est la suivante : À une \struck{\ill{}} \uncertain{forme} réelle $P_{\mathbb{R}}$, on associe l'\uncertain{ensemble} de ses pts réels, qui est un « cercle » dans $P_{\mathbb{R}}(\mathbb{C}) = P(\mathbb{C})$ ($P = P_{\mathbb{R}} \otimes_{\mathbb{R}} \mathbb{C}$) et on prend les deux composantes connexes de $P_{\mathbb{R}}(\mathbb{C}) - P_{\mathbb{R}}(\mathbb{R})$ …. Notons que l'inversion p.r.\ au cercle $P_{\mathbb{R}}(\mathbb{R})$ …

\emph{Scholie} On a un foncteur \add{$\mathbb{P}$} de la catégorie des courbes analytiques non sing.\ connexes et simplement connexes, vers la catégorie des droites projectives complexes, et un homomorphisme fonctoriel
\[
X \xhookrightarrow{\ i_X\ } \mathbb{P}(X)(\mathbb{C})
\]
faisant de $X$ un domaine homogène de $\mathbb{P}(X)$. Un tel foncteur est unique (à isom.\ unique près \uncertain{respectant} $i : \mathrm{id} \to \mathbb{P}(-)(\mathbb{C})$).

Soit maintenant $X$ une courbe analytique non singulière. Pour \uncertain{tout} pt $x \in X$, soit $\widetilde{X}(x)$ le rev.\ universel de $X$ en $x$. C'est un espace analytique simplement connexe, et pour $x$ variable, on trouve un système local (foncteur sur le groupoïde fondamental $\Pi(X)$ de $X$) à valeurs les $\widetilde{X}(x)$ (considérés comme objets de la cat.\ des \add{courbes} \struck{esp.} analytiques connexes et simplement connexes — mais pas pointées). Le système des projections
\[
\widetilde{X}(x) \longrightarrow X = X(x)
\]
est un tel système local d'espaces analytiques. On lui associe donc trivialement un \uncertain{homomorphisme} d'espaces analytiques sur $X$

\begin{tikzcd}
\widetilde{X} \arrow[r] \arrow[dr] & X \times X \arrow[d, "\mathrm{pr}_2"] \\
 & X
\end{tikzcd}

munis de connexions intégrables (en un sens \uncertain{strict} : il y a transport parallèle …). L'hom.\ est compatible avec les connexions, et transforme la section unique $\sigma$ : $x \mapsto \tilde{x}$ de $\widetilde{X}$ en la section diagonale de $X$ (d'où on conclut que $\sigma$ est analytique).

Cette construction marche pour \uncertain{des} \ill{} analytiques. Mais utilisant le fait que $X$ est une courbe non sing., donc les $\widetilde{X}(x)$ sont des courbes non sing.\ connexes et simplement connexes, on trouve que le système local des $\widetilde{X}(x)$ donne naissance au système local des $\mathbb{P}(\widetilde{X}(x))(\mathbb{C})$, d'où un \emph{fibré projectif canonique} sur $X$, noté

\page{127}
$\mathbb{P}_X$, avec une \struck{\ill{}} \add{connexion} intégrable, et une immersion \add{de $X$-esp.\ an.} (respectant les conn.)

\begin{tikzcd}
\widetilde{X} \subset \mathbb{P}_X \arrow[d] \\
X
\end{tikzcd}

Comme $\widetilde{X}$ est muni d'une section unique, $\mathbb{P}_X$ a une section unique analytique. Donc

\emph{Théorème 2} Pour $X$ courbe analytique \add{non sing.} \uncertain{connexe}, on trouve, de façon unique à iso unique près :
\begin{itemize}
\item[a)] Un fibré analytique $\mathbb{P}_X$ sur $X$ en droites projectives
\item[b)] Une connexion (intégrable) dans $\mathbb{P}_X$ (rel : $X/\mathbb{C}$)
\item[c)] Une section $\sigma$ de $\mathbb{P}_X$ sur $X$
\item[d)] Un ouvert $\widetilde{X}$ de $\mathbb{P}_X$ contenant $\sigma(X)$
\item[e)] Un hom.\ de $X$-espaces analytiques $\widetilde{X} \xrightarrow{\ p\ } X_X$ ($= X \times X$)
\end{itemize}
tels que
\begin{itemize}
\item[1°)] $p \circ \sigma = \Delta_X$ (section diagonale)
\item[2°)] $\forall x \in X$, $p_x : \widetilde{X}_x \longrightarrow (X_X)_x = X$ fait de $\widetilde{X}_x$ un rev.\ universel de $X$ (i.e.\ c'est un rev.\ de $X$, \add{la comp.\ connexe de $x$ dans} connexe et simplement connexe)
\item[3°)] $p$ est compatible avec les connexions
\item[4°)] $\forall x \in X$, $\widetilde{X}_x \subset (\mathbb{P}_X)_x$ est un domaine homogène de $(\mathbb{P}_X)_x$.
\end{itemize}

\marginal{\ill{} que une $\widetilde{F}$ d'hom.\ \ill{} \ill{} le long de $\sigma$ \ill{} \ill{} st.\ \ill{} \ill{} le long de $\sigma$) et \ill{} \ill{} analytiques \ill{} plus \ill{}}
\marginal{\ill{} b) et c) (\ill{} que \ill{} pour \ill{} sous \ill{} … ) cf.\ \ill{}}
\marginal{d) et e) \ill{} les faisant \ill{} \ill{} des \ill{} \ill{}}
\marginal{Inversement, 1°), 2°), 3°) \ill{} d) et e) \ill{}}
\note{quatre notes marginales, écrites verticalement le long du bord gauche et reliées par des accolades aux données a) à e) et aux conditions 1°) à 3°) ; elles semblent dire en quoi les données se déterminent l'une l'autre, mais ne se lisent pas}

\page{128}
Supposons $X$ connexe et distinguons \uncertain{ces} 3 cas, suivant le type d'isomorphie de son revêt.\ universel de $X$

1°) Cas \struck{$S^{2} =$} $\mathbb{P}^{1}(\mathbb{C})$. Alors $X$ \struck{$\simeq$} $\mathbb{P}^{1}_{\mathbb{C}}$ est \uncertain{une} droite projective (car \uncertain{simpl.\ connexe}), \uncertain{tout} autour analytique — et en deux \ill{} \ill{} — \ill{} \ill{} connexion triviale, \ill{} \ill{} …), $\widetilde{X} \simeq X \times X$ \add{$= \mathbb{P}^{1}_{X}$ \ill{}}, $p$ est l'identité.

2°) Cas $\mathbb{E}^{1}(\mathbb{C}) = \mathbb{C}$. Alors $X$ est soit \struck{isomorphe} \uncertain{algébrique} (car $X$ \uncertain{comp.\ \ill{}}), soit isomorphe à $\mathbb{C}^{*}$ (cela provient \uncertain{de ce} que les \add{\uncertain{autres} sous-gpes \uncertain{fixes} de $\mathbb{R}^{2}$,} \struck{\ill{} groupes} \struck{discrets} \ill{} de $\mathbb{E}^{1}_{\mathbb{C}}$, \uncertain{c'est kif-kif}) \struck{$\operatorname{Aff}(1, \mathbb{C}) \simeq \mathbb{C}^{*} . \mathbb{C}$ (\ill{} \ill{} \ill{})} \struck{des translations, et les \ill{} groupes discrets \ill{}} \add{sous-groupes discrets} \uncertain{du} groupe des translations $\mathbb{C}$ sont de rang $0$, $1$ (quotient $\mathbb{C}^{*}$) ou $2$ (quotient courbe elliptique). On note que $X$ est muni d'une structure algébrique (tout autant analytique \uncertain{d'une} $\mathbb{E}^{1}_{\mathbb{C}}$, — \uncertain{de} courbe elliptique, — \uncertain{étant} algébriques). Ici $\widetilde{X} = \mathbb{P}_X - \sigma'$,

où $\sigma'$ est une \emph{sect.\ horizontale} de $\mathbb{P}_X$ ; donc, \struck{$\subset \mathbb{P}_X$ \ill{} \ill{}} \uncertain{muni} des deux sections analytiques $\sigma$, $\sigma'$ disjointes, en considérant $\mathbb{P}_X$ comme clôture projective d'un fibré vectoriel analytique $V$, qui \struck{n'est autre} \add{\uncertain{munit} $\widetilde{X}$ comme esp.\ analytique sur $X$} \struck{\ill{} \ill{}} \add{d'une structure vectorielle}. On observe que la connexion de $X$ \struck{\ill{}} n'\uncertain{induit} pas la structure \uncertain{vectorielle}. Mais on peut

\marginal{La cat.\ (pour iso) des \ill{} \ill{} \ill{} \ill{} (\ill{} $\ill{}$) \ill{} \ill{} des \ill{} \ill{} \ill{}}
\note{note marginale écrite en oblique dans l'angle inférieur gauche, reliée par une accolade au cas 2°) ; elle ne se lit pas}

\page{129}
écrire $V = \mathbb{V}(\omega_\sigma)$, $\omega_\sigma$ étant le fibré conormal le long de la section $\sigma$, qui est aussi isom.\ au fibré conormal de $X_X$ le long de la section diagonale, donc
\[
V = \mathbb{V}(\underbrace{\Omega^{1}_{X/\mathbb{C}}}_{\omega_{X/\mathbb{C}}}) = \mathbb{W}(\underbrace{t_{X/\mathbb{C}}}_{\text{fibré tangent}})
\]
Donc $\widetilde{X}$ s'identifie au fibré tangent à $X$, $\sigma$ à la section nulle, $\mathbb{P}_X$ à la clôture projective de $\widetilde{X} = \mathbb{W}(t_{X/\mathbb{C}})$. La connexion \struck{intégrable} \add{plate} canonique de $\mathbb{P}_X$ est associée à une connexion de $\widetilde{X}$ \emph{considéré comme fibré en droites affines}, i.e.\ une connexion sur le fibré principal associé de groupe $\operatorname{Aff}(1, \mathbb{C})$ (contenant $\mathbb{G}_m(\mathbb{C})$ dans $\widetilde{X} = \mathbb{W}(t_{X/\mathbb{C}})$ \uncertain{provient}, et contenu dans $GP(1, \mathbb{C})$).

Il \struck{y aurait lieu} \add{\uncertain{conviendrait}} de \uncertain{caractériser} cette connexion, et prouver (en l'\uncertain{espérant} !) qu'elle est en fait algébrique. \uncertain{Cette} question ne se pose que si $X$ \uncertain{est} compacte, i.e.\ $X = \mathbb{E}^{1}_{\mathbb{C}}$ ou $X = \mathbb{G}_{m\,\mathbb{C}}$ (\uncertain{le cas compact} \emph{réglé} par \emph{GAGA} !). Or dans le cas $X \simeq \mathbb{E}^{1}_{\mathbb{C}}$, on a encore $\widetilde{X} = X \times X \subset \mathbb{P}_X = \mathbb{P}(X) \times X$, \uncertain{puisque} $\widetilde{X} \simeq X$, $p : \widetilde{X} \to X_X$ est l'identité (donc $\sigma = \Delta_X$), tout est clair (et \uncertain{tout} est algébrique — la connexion étant d'ailleurs \uncertain{singulière} à l'infini !)

Dans le cas où $X \simeq \mathbb{C}^{*} \simeq \mathbb{G}_{m}(\mathbb{C})$

\marginal{\uncertain{oui}, elle \uncertain{est} algébrique \ill{} \ill{} \ill{}}
\note{note marginale à gauche, reliée par une flèche à la phrase « qu'elle est en fait algébrique »}

\page{130}
Dans les éléments de structure canoniques associés à la courbe \add{$\mathbb{C}$-}analytique $X$, nous allons retenir maintenant ceux qui « gardent un sens algébrique » lorsque $X$ « est » algébrique — et a fortiori dans le cas propre, grâce à Gaga. Ce sont
\begin{itemize}
\item[a)] Fibré en droites projectives $\mathbb{P}_X$ sur $X$
\item[b)] connexion (plate) $c_X$ sur $\mathbb{P}_X$ (rel.\ $X/\mathbb{C}$)
\item[c)] section $\sigma$ de $\mathbb{P}_X$
\end{itemize}
[d) Un isom.\ de « solutions formelles sur $X$ »
\[
\widehat{\mathbb{P}}_X \xrightarrow[\ \sim\ ]{\ \hat{p}_X\ } \widehat{X}_X
\]
(complétions formelles le long des sections uniques $\sigma$ et $\Delta_X$) respectant les connexions.

D'ailleurs, l'élément de structure d) \uncertain{est contenu} dans b) et d), moyennant la condition (\uncertain{vérifiée} \struck{\ill{}} \add{\uncertain{si} $\exists$ (\ill{} d) \uncertain{compatibles}}) \struck{\ill{} \ill{}} entre \uncertain{eux}) :

1°) \emph{La section $\sigma$ est nulle part stationnaire} (pour la connexion donnée \ill{} sur $\mathbb{P}_X$ …]
\note{« contenu dans b) et d) » se lit ainsi ; on attendrait « b) et c) »}

Même si $X$ n'est pas propre, \struck{\ill{}} les travaux de Deligne montrent l'existence et l'unicité de structures \emph{algébriques} a) b) \add{(unicité, pas l'existence de c) !)} déduites des constructions transcendantes du th.~2, lorsqu'on \struck{fait} \add{\uncertain{rajoute}} de plus la condition :

\struck{\ill{}} La connexion $c_X$ de $\mathbb{P}_X$ est \emph{régulière} à l'infini.

Mais nous allons étudier l'existence (et l'unicité, et préciser ce pb d'existence — pour une autre voie

\page{131}
Quand \uncertain{a-t-on} de telles données \emph{algébriques} sur une courbe algébrique $X$ \add{sur $\mathbb{C}$}, \struck{\ill{}} \uncertain{conduisant} \uncertain{nécessairement} par la relation analytique du th.\ aux termes de celles-ci ? Il \emph{faut} donc construire l'ouvert $\widetilde{X} \subset (\mathbb{P}_X)^{\mathrm{an}}$ et $\widetilde{X} \longrightarrow (X_X)^{\mathrm{an}}$.

Il nous a déjà une bonne \emph{forme} d), et celui-ci pourrait se prolonger, suivant les cas : 1°) en $\mathbb{P}_X \to X_X$ par un \emph{iso} (cas $X$ droite projective) 2°) un \struck{\ill{}} \add{hom.}
\[
\text{\struck{$(\mathbb{P}_X^{\mathrm{an}} \supset \sigma'(X))$}}\quad \widetilde{X} \overset{\mathrm{df}}{=} \mathbb{P}_X^{\mathrm{an}} - \sigma'^{\mathrm{an}}(X^{\mathrm{an}}) \xrightarrow{\ p\ } (X_X)^{\mathrm{an}},
\]
où $\sigma'$ est une section horizontale (\uncertain{univoquement} déterminée) de $\mathbb{P}_X$ sur $X$, faisant de $\widetilde{X}$ un rev.\ de $(X_X)^{\mathrm{an}}$, et transformant $\sigma$ en $\Delta_X$. (\struck{\ill{}} \add{Ici} $p$ est un iso \struck{\ill{}} $X \simeq \mathbb{E}^{1}_{\mathbb{C}}$ i.e.\ $X$ simplement connexe non propre). Le cas où $X \simeq \mathbb{E}^{1}_{\mathbb{C}}$, $\mathbb{G}_{m\,\mathbb{C}}$, ou $X$ courbe elliptique. \emph{Question} A-t-on une description \struck{purement} algébrique de $\sigma'$ (en supposant qu'on ait aux \uncertain{données} de $(\mathbb{P}_X, c_X, p_X)$) ? Et de $p_X$ dans le cas où c'est un iso analytique, i.e.\ $X \simeq \mathbb{E}^{1}_{\mathbb{C}}$ — alors $p_X$ est lui-même algébrique ! 3°) On a un ouvert $\widetilde{X}$ \struck{\ill{}} (\add{dont les fibres sont des} \add{\uncertain{disques}}) \struck{\ill{}} \uncertain{dans} les fibres de $\mathbb{P}_X^{\mathrm{an}}$, lesquels disques sont les domaines d'existence maximaux des $\hat{p}_x : (\mathbb{P}_x)^{\mathrm{an}} \longrightarrow X^{\mathrm{an}}$ …

\marginal{N.B.\ P.~Deligne \ill{} \ill{} \ill{} \ill{} \ill{} \ill{}}
\note{note marginale à gauche, cerclée et barrée d'une croix}
\marginal{\uncertain{On a toujours} \ill{} \ill{} \ill{} $\mathbb{C}$, \ill{} \ill{} \ill{} \ill{} \ill{} \ill{} \uncertain{clôture projective} \ill{} \uncertain{fibré} \ill{} \ill{}}
\note{seconde note marginale, écrite en oblique à gauche des cas 2°) et 3°) ; à peu près illisible}

\page{132}
Or, pour une courbe lisse donnée $X$ sur un schéma de base $S$ \emph{de car.\ nulle} (p.\ ex.\ $S = \operatorname{Spec}\mathbb{C}$) on a donné une description purement alg.\ des structures du type a) b) c) (avec connexion $c_X$ sur $P$ rel.\ à $X/S$), satisfaisant la condition 1° — pour le moment une condition par la condition 2°. La cat.\ de ces structures est rigide — de façon précise, les $P$ (fibré de S.D.\ de rang 1 sur $X$) sont can.\ isom.\ à $\mathbb{P}(E)$, où $E$ est l'extension
\[
0 \longrightarrow \underline{\omega}_{X/S} \longrightarrow E \longrightarrow \underline{\mathcal{O}}_X \longrightarrow 0
\]
\note{sous $\underline{\omega}_{X/S}$, un signe $\Vert$ le renvoie à $\underline{\omega}$}

donnée par
\[
E \simeq \mathcal{J}/\mathcal{J}^{3} \otimes \omega^{-1}
\]
($\mathcal{J} \subset \mathcal{P}^{\infty}_{X/S}$ l'idéal d'augmentation, donc $\mathcal{J}/\mathcal{J}^{2} \simeq \omega_{X/S} = \omega$), la section $\sigma$ est alors définie par cette structure d'extension, enfin la connexion $c$ sur $P = \mathbb{P}(E)$ est déduite d'un isom.\ $p : \mathcal{P}^{\infty}_{X/S} \xrightarrow{\ \sim\ } \underline{\mathcal{O}}_{\hat{P}}$ ($\hat{P}$ complété formel de $P$ le long de $\sigma$), prolongeant l'isom.\ $\add{p(\sigma)}$ (et l'ind.\ 2 \uncertain{canonique}) impliqué par la description de $P = \mathbb{P}(E)$ par le choix particulier de $E$. On a un produit, par les isom.\ $\hat{p}$ tels que la conn.\ de $\hat{p}$ déduite par transport de structure de celle de $\mathcal{P}^{\infty}_{X/S}$ (\uncertain{provenant} de la conn.\ tautologique de \struck{$X \times X$} $X \times_S X$ sur $X$) \uncertain{provienne} \uncertain{bien} d'une connexion \emph{de $P$}. On constate donc que ces iso (correspondant
\note{la phrase se poursuit à la page suivante}

\page{133}
donc à toutes les structures possibles du type envisagé sur $X$) sont en correspondance biunivoque, i.e.\ \uncertain{vérifient}, avec l'ens.\ des relèvements $p_{(3)}$ de $p_{(2)}$ à l'ordre~3 (sans autre condition sur ceux-ci). Ce sont \add{donc} les sections sur $X$ d'un torseur canonique sur $S$, sous le faisceau abélien $\underline{\mathrm{Hom}}(\underline{\omega}, \underline{\omega}^{\otimes 3}) \simeq \underline{\omega}^{\otimes 2}$ :
\[
\text{\struck{$\mathcal{P}$}} \quad \mathcal{C}_{X/S} \quad \text{torseur sous } \underline{\omega}^{\otimes 2} = \underline{\Omega}^{1\,\otimes 2}_{X/S}
\]
\note{sous $\mathcal{C}_{X/S}$ : « \uncertain{ensemble} des “connexions projectives” »}

On dit aussi que les structures du type envisagé sur $X$ sont des \emph{connexions projectives} sur $X$, et le torseur envisagé est le \emph{torseur des connexions projectives} sur $X/S$. [Elles sont aussi en correspondance biunivoque avec l'ens.\ des connexions \add{\uncertain{correspondantes}} $c$ sur $P = \mathbb{P}(E)$ ($E$ comme dessus), ou ce qui revient au même, des connexions \add{(\uncertain{compatibles} \ill{} \ill{})} sur le fibré en $M_3$ \struck{\ill{}} \uncertain{lié} associé (i.e.\ qui \add{\uncertain{dirait}} \struck{\ill{}} $P \to X$ de $\mathcal{J}_{R/X}$, \uncertain{on} \uncertain{aura} $E \otimes \check{E}/\underline{\mathcal{O}}_X$) — on trouve
\[
\mathcal{V} \simeq (\operatorname{Sym}^{2} E) \otimes \underline{\omega}^{-1}
\]
extension successive — \uncertain{commençant} \ill{} \ill{} — de $\underline{\omega}^{1}, \underline{\omega}^{0}, \underline{\omega}^{-1}$ — et il faut prendre les connexions $\nabla$
\[
c : \mathcal{V} \longrightarrow \mathcal{V} \otimes \underline{\omega} \simeq \operatorname{Sym}^{2} E
\]
qui, sur $\underline{\omega} \subset \mathcal{V}$ (\uncertain{défini} nécessairement dans $\operatorname{Ker}(\operatorname{Sym}^{2} E \to \operatorname{Sym}^{2}\underline{\mathcal{O}}_X \simeq \underline{\mathcal{O}}_X)$, ext.\ de $\underline{\omega}$ par $\underline{\omega}^{2}$,
\note{la phrase se poursuit à la page suivante}

\page{134}
laquelle extension s'identifie à $\mathcal{P}^{1}_{X/S}(\omega)$ par sa structure d'extension naturelle de $\omega$ par $\underline{\omega} \otimes \underline{\Omega}^{1}_{X/S} \simeq \underline{\omega} \otimes \underline{\omega} = \underline{\omega}^{2}$) s'identifie à l'op.\ diff.\ de degré~1 (rel.\ à $S$) universel sur $\underline{\omega}$
\[
c|\underline{\omega} = d^{1}_{\underline{\omega}, X/S} : \underline{\omega} \longrightarrow \mathcal{P}^{1}_{X/S}(\underline{\omega}) .
\]
L'indéterminati\supplied{on} \add{dans \uncertain{cette}}, \uncertain{sur} cette condition, au lieu d'être dans $\Gamma(\underline{\omega} \otimes \mathcal{V})$ — \struck{\ill{}} (\uncertain{connexion} sur $\mathcal{V}$ \add{\uncertain{par} \uncertain{produit}, i.e.\ \uncertain{dans}} \add{une section} des $\mathcal{V} \otimes \underline{\omega}$) est dans $(\underline{\omega} \otimes \underline{\mathrm{Bil}}_0 \simeq \underline{\omega} \otimes \underline{\omega} \simeq \underline{\omega}^{\otimes 2}$ — on retombe sur ses pattes !

Ceci « rappelé », on voit donc que si $S = \operatorname{Spec}\mathbb{C}$, donc $X$ courbe algébrique lisse$/\mathbb{C}$, \uncertain{dans} la \struck{\ill{}} \add{\uncertain{construction}} \emph{transcendante} \struck{\ill{}} \add{\uncertain{indiquée}} plus haut \struck{\ill{} \ill{} des} \add{s'\uncertain{interprète} comme la} construction d'une connexion projective (canonique) sur $X/\mathbb{C}$, i.e.\ d'une section canonique du torseur $\mathcal{C}_{X/\mathbb{C}}$ (des connexions projectives sur $X$), \add{sous $\underline{\omega}^{\otimes 2}_{X/\mathbb{C}} = \underline{\omega}^{2}$} \add{on se demande quand} (cette connexion projective est de plus « régulière » à l'infini » i.e.\ la connexion correspondante sur $\mathbb{P}_X \simeq \mathbb{P}(\mathcal{J}/\mathcal{J}^{3})$ ($\mathcal{J}$ idéal d'augmentation de $\mathcal{P}^{\infty}_{X/\mathbb{C}}$) \struck{étant} \add{est} régulière à l'infini. Sauf erreur, cela doit signifier

\marginal{\ill{} \ill{} \ill{} \ill{} \ill{} \ill{} \ill{} \uncertain{projectives} \ill{} \ill{} \uncertain{régulière} \ill{} \ill{} \ill{} \uncertain{donnée} \ill{} \ill{}}
\marginal{\ill{} \ill{} $\mathbb{E}^{1}_{\mathbb{C}}$ \ill{} $\mathbb{G}_{m\,\mathbb{C}}$ \ill{} $\mathbb{P}$ \ill{} \ill{} \ill{} \ill{} \ill{} \ill{}}
\note{deux notes marginales écrites verticalement le long du bord gauche, la seconde reliée par un trait courbe à la condition « régulière à l'infini » ; à peu près illisibles}

\page{135}
que cette section du torseur $\mathcal{C}$ (qui est en fait défini sur $\hat{X}$, compactifiée non singulière \uncertain{naturelle} de $X$,) sont $\mathcal{C}_{\hat{X}/\mathbb{C}}$) est \uncertain{admet} « à l'infini » $\hat{X} - X$ que des pôles simples au plus, i.e.\ si $D$ est le diviseur correspondant \add{sur $\hat{X}$} à multiplicités toutes 1, de support $\hat{X} - X$, on aura comme groupe d'indéterminations des conn.\ proj.\ rég.\ à l'infini sur $\hat{X}$
\[
H^{0}(\hat{X}, \underline{\omega}^{\otimes 2}_{\hat{X}}(D))
\]
(qu'on peut relier aussi aux invariants classiques de la courbe complétée $\hat{X}$ les \uncertain{écrivant} \struck{\uncertain{\ill{}}} d'abord $\underline{\lambda} = $ \struck{$\underline{\omega}^{2}$} \ldots{} \uncertain{lieu} de $\underline{\omega}^{\otimes 2}_{\hat{X}/\mathbb{C}}$)
\[
\text{\struck{$0 \to \underline{\lambda} \to \underline{\lambda}(D) \to$}} \quad \underline{\omega}(D)/\underline{\omega}^{2} \longrightarrow 0
\]
\[
0 \longrightarrow \underline{\lambda} \longrightarrow \underline{\lambda}(D) \longrightarrow \underline{\lambda}(D)/\underline{\lambda} \longrightarrow 0
\]
déduite par $\otimes\,\underline{\lambda}$ de
\[
0 \longrightarrow \underline{\mathcal{O}}_{\hat{X}} \longrightarrow \underline{\mathcal{O}}_{\hat{X}}(D) \longrightarrow \underline{\mathcal{O}}_{\hat{X}}(D)/\underline{\mathcal{O}}_{\hat{X}} \longrightarrow 0
\]
où $\underline{\mathcal{O}}_{\hat{X}}(D) = \mathcal{J}_D^{-1}$ (idéal de $D$) et $\underline{\mathcal{O}}_{\hat{X}}(D)/\underline{\mathcal{O}}_{\hat{X}} \simeq (\mathcal{J}_D/\mathcal{J}_D^{2})^{-1} \simeq$ \struck{\ill{}} $(\underline{\omega}^{-1}|D)$, d'où
\[
\underline{\lambda}(D)/\underline{\lambda} \simeq (\underline{\lambda}\,\underline{\omega}^{-1})|D
\]
et la suite exacte
\[
0 \to H^{0}(\hat{X}, \underline{\lambda}) \to H^{0}(\hat{X}, \underline{\lambda}(D)) \to H^{0}(D, \underline{\lambda}\,\underline{\omega}^{-1}|D) \to H^{1}(\hat{X}, \underline{\lambda})
\]
(ici $\underline{\lambda} = \underline{\omega}^{\otimes 2}_{\hat{X}/\mathbb{C}} \simeq \underline{\omega}^{2}$).

\marginal{On suppose $X$ connexe}
Si $\hat{X}$ est de genre $g \geqslant 2$, $H^{1}(\hat{X}, \underline{\lambda}) = 0$ (i.e.\ le dual de $H^{0}(\hat{X}, \underline{\omega}^{-1}) = 0$) donc
\[
0 \to H^{0}(\hat{X}, \underline{\omega}^{2}) \to H^{0}(X, \underline{\omega}^{2})^{\mathrm{reg}} \to H^{0}(D, \underline{\omega}|D) \to 0
\]
\note{sous les termes, de gauche à droite : $H^{0}(\hat{X}, \underline{\omega}^{2}) \simeq \underline{\Omega}^{1}_{M_g/\mathbb{C}}(s)$, de rang $3g-3$ (cerclé), $M_g$ variété modulaire des courbes de genre $g$ \uncertain{non compactes} ; $H^{0}(X, \underline{\omega}^{2})^{\mathrm{reg}} \simeq$ (\uncertain{can.} ?) $\underline{\Omega}^{1}_{M_{g,\nu}/\mathbb{C}}(s)$, $M_{g,\nu}$ variété modulaire des courbes de genre $g$ à $\nu$ trous ; $H^{0}(D, \underline{\omega}|D)$ de rang $\nu = \operatorname{card} D$}

\page{136}
(b) $g \leqslant 1$, i.e.\ $g = 0, 1$. Est-ce que
\[
H^{1}(\hat{X}, \underline{\lambda}) \longrightarrow H^{1}(\hat{X}, \underline{\lambda}(D)) \quad \text{est inj.\ ?}
\]
On a encore
\[
H^{0}(\hat{X}, \underline{\omega}^{-1}) \longleftarrow H^{0}(\hat{X}, \underline{\omega}^{-1}(-D))
\]
\note{sous $\underline{\omega}^{-1}$ : $= \underline{t}_{\hat{X}/\mathbb{C}}$, fibré tangent, et $H^{0}(\hat{X}, \underline{\omega}^{-1}) = \operatorname{Lie}\underline{\mathrm{Aut}}(\hat{X})$ (cerclé) ; sous $\underline{\omega}^{-1}(-D)$ : $= \underline{t}_{\hat{X}/\mathbb{C}}(-D)$, champs tangents s'annulant sur $D$, et $H^{0} = \operatorname{Lie}\underline{\mathrm{Aut}}(X)$ (cerclé)}

\struck{injectif}, surjectif ? Si son conoyau \add{sans pôles} $\check{Q}(X)$, on aura
\[
0 \to H^{0}(\hat{X}, \underline{\omega}^{2}) \to H^{0}(X, \underline{\omega}^{2})^{\mathrm{rég}} \to H^{0}(D, \underline{\omega}|D) \to \check{Q}(X) \to 0
\]
\emph{Cas $g = 1$}, \add{\uncertain{dans}} dimension $H^{0}(\hat{X}, \underline{t}_{\hat{X}/\mathbb{C}}) =$ \struck{\ill{}} 1 et si $D \neq 0$ i.e.\ $X \neq \hat{X}$ i.e.\ $X$ non complète, on a $Q(X) \simeq H^{0}(\hat{X}, \underline{\omega}^{-1}) \simeq t_{\hat{X}}$ (\add{\uncertain{algèbre de Lie}} \struck{\ill{}} du groupe \uncertain{des translations}) d'où $\check{Q}(X) = \omega_{\hat{X}} = \check{t}_{\hat{X}}$ (formes diff.\ inv.\ sur \uncertain{\ill{}}). Interprétation plausible \uncertain{sur} \uncertain{rigidité} du pb modulaire des courbes \emph{elliptiques} (i.e.\ \emph{pt marqué}) avec $\nu$ pts \add{\uncertain{distincts}} \emph{\uncertain{de somme nulle}} \struck{\uncertain{\ill{}}} \add{(\uncertain{cas} \uncertain{où} \uncertain{\ill{}} \uncertain{à} \ill{})} — courbes elliptiques \uncertain{tenues} \uncertain{à} \ill{} et $\nu$ pts \add{(non ordonnés)} \struck{marqués} distincts et un pt choisi \uncertain{à} \uncertain{leur} \uncertain{somme} … \struck{\ill{}}
\[
\text{Alors } H^{0}(\hat{X}, \underline{t}_{\hat{X}/\mathbb{C}}) = 0
\]
\emph{Cas $g = 0$}. $H^{0}(\hat{X}, \underline{t}_{\hat{X}/\mathbb{C}})$ est de dim~3, et suivant les cas
\[
\begin{array}{ll}
\nu = 1 & \dim H^{0}(\hat{X}, \underline{t}_{\hat{X}/\mathbb{C}}(-D)) = 2\\
\nu = 2 & \phantom{\dim H^{0}(\hat{X}, \underline{t}(-D))} = 1\\
\nu \geqslant 3 & \phantom{\dim H^{0}(\hat{X}, \underline{t}(-D))} = 0
\end{array}
\]
donc $Q(X)$ est respectivement de dim $1$, $2$ ou $3$.

Donc pour $0 \leqslant \nu \leqslant 3$, on aura $\dim \check{Q}(X) = \operatorname{card}\operatorname{supp} D = \dim H^{0}(D, \underline{\omega}|D)$
\note{la ligne « Alors $H^{0}(\hat{X}, \underline{t}_{\hat{X}/\mathbb{C}}) = 0$ » est reliée par une accolade au cas $g = 0$ qui suit ; la conclusion, « $0 \leqslant \nu \leqslant 3$ », se lit ainsi alors que le cas $\nu \geqslant 3$ vient de donner $\dim Q(X) = 3$ : on transcrit tel quel}

\page{137}
donc $H^{0}(X, \underline{\omega}^{2})^{\mathrm{rég}} = 0$ : il n'y a qu'une \emph{seule connexion projective} \add{régulière à l'$\infty$ !} sur $X$ si \emph{$X$ est de genre $0$ et à $\nu \leqslant 3$ trous}. Par contre pour $\nu \geqslant 4$, l'espace $H^{0}(X, \underline{\omega}^{2})^{\mathrm{rég}}$ est de dim $\nu - 3$ — c'est encore la dim de l'espace modulaire pour la droite projective à $\nu$ trous, et sans doute on trouve un isom.\ can.
\[
H^{0}(X, \underline{\omega}^{2})^{\mathrm{rég}} \simeq \underline{\Omega}^{1}_{M_{0,\nu}}(s) \quad \ldots
\]
Bien entendu, dans le cas $\nu \leqslant 3$, la seule et unique connexion projective sur $X$ est celle induite par la connexion projective de $\hat{X}$ (et \uncertain{\ill{}} \uncertain{toutes} \uncertain{celles} de \uncertain{\ill{}} \uncertain{transcendant} !) Dans le cas de $\nu \leqslant 3$, \add{\ill{}} \add{\uncertain{possibilités}} \uncertain{dans} \uncertain{sa} \uncertain{description} \uncertain{par} \uncertain{tout} \uncertain{\ill{}} \uncertain{transcendant} ! \struck{\ill{}} la connexion projective \add{\uncertain{unique}} (\struck{\uncertain{\ill{}}} \add{\uncertain{\ill{}}}) \uncertain{de} $\mathbb{P}_X = \hat{X} \times X$, munie de la section diagonale (et de la connexion triviale. Dans le cas de $X$ quelconque (de genre $0$) cette connexion est \uncertain{après} définition, et fournit une \emph{section canonique}
\marginal{\ill{} \ill{} \ill{} \ill{} \ill{} \ill{} \ill{} \ill{}}
\note{note marginale écrite en oblique à gauche, à la hauteur de « transcendant ! » ; illisible. La phrase finale reste en suspens au bas de la page, dont le reste est blanc}

\page{138}
\add{On suppose \ill{} inv.\ sur $S$.}

\emph{Prop.}\ Soit $H$ \uncertain{schéma} \uncertain{en} groupes lisse sur $S$, de dim rel.~1, et $X$ un torseur sous $H$. Soit $(\mathcal{V}, \mathcal{L}, \underline{\omega})$ \uncertain{l'alg.\ de Lie} de rang \uncertain{lin.}~3 \uncertain{correspondante} \add{$(\mathbb{P}, \sigma)$, le fibré proj.\ : section associée} sur $X$. Il y \uncertain{a} \uncertain{un} \uncertain{fibré} par translation et \uncertain{le} \uncertain{descend} par
\[
\underset{\omega_0}{\underline{\omega}_0} \subset \underset{\underline{\mathcal{O}}_S}{\mathcal{L}_0} \subset \underset{\omega_0^{-1}\ \text{sur } X}{\mathcal{V}_0}, \qquad (P_0, \sigma_0) \quad P_0 = \mathbb{P}(\mathcal{L}_0)
\]
\[
\mathcal{L}_0 = (\underline{\mathcal{O}}_H/\mathcal{J}^{3})^{+}, \quad \mathcal{J} \text{ id.\ d'augm.\ de } H
\]
Alors les conn.\ \uncertain{projectives} \add{\uncertain{invariantes} sous $H$} sont en corresp.\ 1--1 avec \add{les sections d'}un \uncertain{certain} sous-torseur \struck{\ill{}} de $\mathcal{V}_0 \otimes \underline{\omega}_0$ sous $\underline{\omega}_0 \otimes \underline{\omega}_0 = \underline{\omega}_0^{2}$, se projetant sur la section $1$ de $\underline{\omega}_0^{-1} \otimes \underline{\omega}_0 \simeq \underline{\mathcal{O}}_S$, [\uncertain{formée} des sections \add{\uncertain{locales}} \struck{\ill{}} de $\mathcal{V}_0 \otimes \underline{\omega}_0$ qui \uncertain{sont} \uncertain{compatibles} \uncertain{à} $\mathcal{L}_0 \otimes \underline{\omega}_0$, \add{\uncertain{suivant} celle de Hausdorff}] \uncertain{on} \uncertain{a} \struck{\ill{}} \uncertain{l'ensemble} \uncertain{canonique} \uncertain{de}
\[
\mathcal{V}_0 \simeq \underline{\omega}_0 \oplus \underline{\mathcal{O}}_S \oplus \underline{\omega}_0^{-1}
\]
\uncertain{\ill{}} \uncertain{\ill{}} \uncertain{suivant} sous la forme
\[
\omega_2 + \omega_1 + 1 \in \Gamma(\underline{\omega}_0^{2} + \underline{\omega}_0 + \underline{\mathcal{O}}_S).
\]
Les sections de ce sous-torseur correspondent aux sections de $\mathbb{P}(\omega_0 + \omega_0^{-1})$ qui ne \uncertain{disparaissent} \uncertain{\ill{}} de la section canonique \uncertain{déduite} de $\omega_0 \subset \omega_0 \oplus \omega_0^{-1}$, — \uncertain{\ill{}} \uncertain{\ill{}} aux \uncertain{scindages} de $0 \to \omega_0 \to \omega_0 \oplus \omega_0^{-1} \to \omega_0^{-1} \to 0$. Il y a \uncertain{donc} une canonique \uncertain{parmi} \uncertain{elles} \uncertain{-ci} \uncertain{celle} de \uncertain{Campbell}-Hausdorff) \uncertain{correspondant} à la section $0 + 0 + 1$ de $\mathcal{V}_0 \otimes \underline{\omega}_0$. Elle est \emph{caractérisée} \add{\uncertain{parmi} les conn.\ \uncertain{proj.}\ inv.} \add{\uncertain{loc.}\ sur $S(\ell)$} par la condition que \struck{\ill{}} l'on \uncertain{puisse} trouver un $P$ une section \add{\uncertain{invariante} par translations} (\uncertain{du} \uncertain{groupe} \uncertain{structural} \ill{}) \uncertain{à} $G_m$, i.e.\ qu'on puisse trouver une

\marginal{\uncertain{\ill{}} \uncertain{\ill{}} \ill{} \ill{} $[\underline{\mathcal{O}}_S] \subset \mathcal{L}_0$ \uncertain{correspond} aux parties d'ordre~2 des fonctions \struck{\ill{}} \uncertain{invariantes} ; on a $[\underline{\mathcal{O}}_S] \otimes \underline{\omega}$, \uncertain{\ill{}} \uncertain{\ill{}} : $\omega_0 \to \mathcal{L}_0 \otimes \omega_0$}
\marginal{\ill{} \ill{} \ill{} \ill{} \ill{} \ill{} $\mathbb{V}(\omega)$ \ill{} \ill{} \ill{}}
\note{deux notes marginales à gauche ; la seconde, écrite verticalement, à peu près illisible}

\page{139}
\uncertain{$G$-}extension sur $X$ invariante par translation
\[
0 \longrightarrow \underline{\omega}_X \longrightarrow F \longrightarrow \underline{\mathcal{O}}_X \longrightarrow 0
\]
et un $G$-iso.\ $P \simeq \mathbb{P}(F)$ ($\simeq \overline{\mathbb{V}(\mathcal{T})}$, où $\mathcal{T}$ est le $G$-torseur \uncertain{associé} sur $\underline{\mathcal{O}}_X$, avec \uncertain{la} \uncertain{structure} \uncertain{tautologique} de $\underline{\omega}_X$ définie par l'extension …) [La \add{\uncertain{trace} $\sigma'$} section$/$de $P$ qui y correspond est \uncertain{toute} \uncertain{disjointe} de $\sigma$, donc correspond à un \uncertain{scindage} \add{(\uncertain{inv.}\ par $H$)} $\mathcal{V} \simeq \underline{\omega} \oplus \underline{\mathcal{O}}_X \oplus \underline{\omega}^{-1}$, \uncertain{qu'on} \uncertain{peut} \uncertain{écrire} \uncertain{par} \uncertain{celui} [\uncertain{que} d'\uncertain{abord} un \uncertain{scindage} \ill{} \ill{} \ill{}] \uncertain{et} \uncertain{d'autre} \uncertain{un} \uncertain{scindage} de Campbell-Hausdorff, \uncertain{sur} $S$ déjà \uncertain{canonique} …) avec $\underline{\mathcal{O}}_X + \underline{\omega}^{-1}$ \add{\uncertain{et} $\underline{\omega}^{-1}$} (\uncertain{invariant} par la connexion), et elle \uncertain{induit} \struck{\uncertain{\ill{}}} \uncertain{\ill{}} \uncertain{triviale} \uncertain{sur} \uncertain{le} [\uncertain{quotient} $\underline{\mathcal{O}}_X$ \ill{}]. Pour les conn.\ \uncertain{projectives} invariantes, \uncertain{pour} \uncertain{lesquelles} on \uncertain{dispose} \uncertain{de} \uncertain{ce} \uncertain{que} l'on peut \uncertain{dire} de la « \uncertain{distinction} », il y a \uncertain{$G$-réduction} (\add{\uncertain{\ill{}}}) \uncertain{du} groupe structural \struck{\ill{}} \add{\uncertain{\ill{}} \uncertain{de}} à un \uncertain{sous-groupe} du $G_0 = GP(\mathbb{P}(\underline{\omega}_0 + \underline{\mathcal{O}}_S))/$. Si on \uncertain{est} \uncertain{sur} un corps alg.\ clos (\uncertain{S spectre} \uncertain{de} $k$), \uncertain{alors} la connexion proj.\ can.\ est \uncertain{caractérisée}, \uncertain{parmi} celles qui sont invariantes, par la condition \struck{\ill{}} que $P$ \uncertain{admette} une section \struck{\uncertain{horizontale}} \add{\uncertain{\ill{}} \uncertain{\ill{}}} (\emph{N.B.}\ \uncertain{en} car.~$0$, \struck{\ill{}} \add{\uncertain{\ill{}}}) \uncertain{\ill{}} \uncertain{invariante}), \uncertain{toutes} les autres en \uncertain{ont} \uncertain{exactement} \uncertain{deux} …

\emph{N.B.}\ Si $H$ est propre sur $S$, \uncertain{i.e.}\ \uncertain{fibrés} connexes (\uncertain{donc} $X/S$ \uncertain{est} \uncertain{schéma} \uncertain{en} \uncertain{courbes} \uncertain{elliptiques}) \uncertain{alors} les connexions \uncertain{projectives} sur $X$ sont \uncertain{toutes} inv.\ par $H$. \uncertain{Ce} qui \uncertain{caractérise} \uncertain{donc} la conn.\ \uncertain{proj.}\ canonique \uncertain{parmi} \uncertain{toutes} les connexions …

\marginal{\ill{} \ill{} \ill{} \ill{} \ill{} \ill{} \ill{} \ill{} \ill{} \ill{}}
\note{le long du bord gauche, sur toute la hauteur de la page, plusieurs colonnes de notes marginales écrites verticalement et reliées par des accolades ; à peu près illisibles. En face de la phrase sur la condition, un « ? » en marge}

\end{document}
